\subsection{Download Malicious Tarballs}
\label{sec:DownloadMaliciousTarballs}

In this subsection, the process and the developed script 
used to retrieve the tarballs of the malicious versions of the packages specified 
in \cref{sec:malwareDataset} are described.

\subsection*{Public Repositories of Malicious Packages}
\label{sec:PublicRepositoriesMaliciousPackages}

The first approach used was to consult public datasets on GitHub, such as
the one by \textcite{OpenSourceDatasetMaliciousSoftwarePackages}.

These public repositories, which contain malicious versions of packages from a registry, 
are active and regularly updated.
Their limitation is that they do not always contain specific malicious versions of compromised packages.

In fact, in our case, the versions of the packages compromised 
by the attacks in \cref{ch:AnalysisRecentAttacks} that we were looking for 
were not initially present in the dataset by \textcite{OpenSourceDatasetMaliciousSoftwarePackages}.

\subsection*{Mirror registry}
\label{sec:MirrorRegistry}

The alternative used was to consult mirror registries, 
which are servers that replicate the content of official package 
repositories such as \ac{NPM}, \ac{PyPI}, and RubyGems \cite{guo2023empiricalstudymaliciouscode}.

Some examples of mirror registries include Huaweicloud, npmmirror, Tencent, and Alibaba.

The main characteristic of a mirror is its synchronization frequency with the official registry.
Such synchronization varies depending on the mirror and can be more or less frequent.
However, since synchronization does not occur instantaneously,
mirrors can preserve versions of packages already removed from the official registry for extended periods.

This latency makes mirrors a useful tool for retrieving malicious versions of packages after their removal from official registries.

\subsection*{Script to Download Malicious Tarballs}
\label{sec:ScriptToDownloadMaliciousTarballs}

To automate the search and download across different mirrors,
the script \path{downloadpkgnpm.py} was developed.

Implemented in Python and available on GitHub\footnote{\url{https://github.com/Frazzerz/Download-malicious-npm}},
the script can be invoked from the command line in the following way:

\begin{verbatim}
python3 downloadpkgnpm.py json_file.json
\end{verbatim}

The script requires as input a \ac{JSON} file where one can specify
a list of mirrors to query and a list of packages and versions to download.

For each package, the script sequentially queries the specified mirrors
in search of the required version.

If the tarball is found, the script proceeds to download it.
Otherwise, if the search fails on all mirrors, the operation's failure is notified.

In both cases, the execution continues automatically with the search for the next package until the list is exhausted.

With this script, it was possible to retrieve the malicious versions 
of the compromised packages used in the tool \path{Analyzer-npm-package-releases} (\cref{sec:ToolForAnalyzingNpmPackages}) 
to create the Malware dataset (\cref{sec:malwareDataset}).